%%%% h2snn_ijcai27_v3b.tex -- IJCAI-27 version of H2S-NN, model v3-B (methodology spec v1.14)
%%%% Experiments are unfinished: every number that does not exist yet is \res.

\documentclass{article}
\pdfpagewidth=8.5in
\pdfpageheight=11in

\usepackage{ijcai27}

\usepackage{times}
\usepackage{soul}
\usepackage{url}
\usepackage[hidelinks]{hyperref}
\usepackage[utf8]{inputenc}
\usepackage[small]{caption}
\usepackage{graphicx}
\usepackage{amsmath}
\usepackage{amssymb}
\usepackage{amsthm}
\usepackage{booktabs}
\usepackage{multirow}
\usepackage[switch]{lineno}

% Comment out this line in the camera-ready submission
\linenumbers

\urlstyle{same}

\newtheorem{proposition}{Proposition}

\newcommand{\R}{\mathbb{R}}
\newcommand{\method}{\textsc{H2S-NN}}
\newcommand{\wisp}{\textsc{WISP}}
\newcommand{\norm}[1]{\left\lVert #1\right\rVert}
\newcommand{\transpose}{\mathsf{T}}
\newcommand{\concat}{\mathbin\Vert}
\newcommand{\res}{{\scriptsize[result]}}
\newcommand{\inc}{\trianglelefteq}

\pdfinfo{
/TemplateVersion (IJCAI.2027.0)
}

\title{H2S-NN: Meta-Path-Structured Sheaf Neural Networks\\ for Heterogeneous Heterophilic Graphs}

% Anonymous submission
\author{
    Anonymous submission
}

\begin{document}

\maketitle

\begin{abstract}
Heterogeneous graphs are frequently heterophilic: entities of one type link to entities of other types whose classes differ, so adjacency does not imply label agreement. Existing methods split the graph into one homogeneous view per meta-path schema and fuse the views; heterophily is then handled on derived graphs, and every node pair joined under the same schema is treated alike, which discards the intermediate entities and the particular meta-path instances that join the pair. We ask whether heterophily can be handled on the original graph itself, by making meta-path instances part of the structure the model operates on. We propose \method{}, a sheaf neural network in which each meta-path instance, sampled by random walks without a meta-path list, is a face of the graph's own sheaf: its restriction maps are forced by the maps of its edges, so no parameters are added, and every node of the instance, including the intermediate ones, is coupled to it. Nodes, edges and faces form one quadratic energy whose minimiser is unique for all parameter values and gives the node representations. On eight node-classification benchmarks and eight link-prediction tasks spanning a range of heterophily, \res{}; controlled variants on identical sampled instances show \res{}.
\end{abstract}

%======================================================================
\section{Introduction}
\label{sec:introduction}

Heterogeneous graphs are a common representation for modeling complex relational structures, such as citation networks \cite{liu2014meta} or social interactions \cite{singh2025overview}. Unlike homogeneous graphs, they contain multiple types of nodes and relations with distinct semantic roles \cite{sun2013mining}. Moreover, many of them are also heterophilic, meaning that adjacent nodes may have different labels or attributes, which makes similarity-based neighbor aggregation unreliable and lead to over-smoothing \cite{lin2025heterophily}. A graph can be both heterogeneous and heterophilic, which means the graphs can have multiple node types while adjacent nodes have different attributes. This makes learning considerably harder than under either property alone. Because an edge joins nodes of different types, which share neither a feature space nor a label space, agreement between neighbors is not even defined and can only be examined between target-type nodes connected through other nodes. Furthermore, the nodes in between decide whether such a pair agrees. In Figure~\ref{fig:example}(a), papers P1 and P9 share the author A7 but belong to different fields, whereas a narrow subject such as S3 may link papers of one field. Methods designed for heterophily usually assume a single node type, and methods designed for heterogeneity usually assume homophily, so neither addresses both properties together.

To handle theMeta-paths are the standard way to handle this situation. A meta-path schema such as \texttt{Paper-Author-Paper} (PAP) describes a composite relation between two target-type nodes that neither node expresses alone \cite{sun2011pathsim}, and a meta-path instance such as P1--A7--P9 is one concrete occurrence of it. Meta-paths connect target nodes whose labels are comparable, distinguish the semantics of different relations (PAP is a shared author, PSP a shared subject), and make heterophily measurable on heterogeneous graphs \cite{lin2025heterophily}. They therefore underlie most heterogeneous graph models \cite{wang2019han,fu2020magnn,yang2023sehgnn}.

\begin{figure*}[t]
  \centering
  \includegraphics[width=0.92\textwidth]{figure/h2snn_example_v3b2.pdf}
  \caption{Running example. \textbf{(a)} A heterogeneous heterophilic citation graph: papers P1, P4, P9 (colour = label), author A7, subject S3. \textbf{(b)} Splitting by meta-path schema builds one paper graph per schema (PAP, PSP); A7 and S3 appear in neither. \textbf{(c)} \method{} keeps the meta-path instance P1--A7--P9 as a face of the original graph: its stalk (dot) is coupled to every node and edge of the instance, including A7, which passes the result on to P4. No edge P1--P9 is added.}
  \label{fig:example}
\end{figure*}

We answer with the \textbf{Heterogeneous Heterophilic Sheaf Neural Network (\method{})}. Following the usual vocabulary of cellular sheaves, we call a node a 0-cell, an edge a 1-cell, and a two-dimensional cell (2-cell) a \emph{face}. \method{} makes each sampled meta-path instance a face of the graph's own sheaf and lets nodes, edges and faces interact through one energy. Our contributions are as follows:
\begin{enumerate}
  \item \textbf{Faces of the graph's own sheaf.} A face is one meta-path instance, never closed into a loop. Its restriction maps are forced by the maps of its edges, adding no parameters, and the face is coupled to every node of the instance (Propositions~\ref{prop:forced} and~\ref{prop:reach}).
  \item \textbf{One energy, one solution.} Node, edge and face stalks interact through a single quadratic energy, whose minimiser is unique for every parameter value (Proposition~\ref{prop:unique}); dropping faces, collapsing them to endpoint edges, or using them as features are configurations of this energy, which makes controlled ablations possible.
  \item \textbf{Schema-free sampling with instance weights.} \wisp{} samples instances by random walks without a meta-path list, with unbiased sampling weights, and a learned weight scores each instance from its own nodes and relations.
\end{enumerate}

%======================================================================
\section{Related Work}
\label{sec:related}

\paragraph{Heterogeneity and heterophily with meta-paths.} HAN \cite{wang2019han} and SeHGNN \cite{yang2023sehgnn} aggregate over one homogeneous view per schema, and LMSPS \cite{li2024lmsps} searches for useful schemas. The heterophily-aware models H2G-former \cite{lin2025heterophily}, Hetero$^2$Net \cite{li2025hetero2net}, H2SGNN \cite{lu2026addressing} and LatGRL \cite{shen2025heterophily} also rely on meta-path-based or latent homogeneous graphs. Intermediate nodes of a schema do not appear in these views. MAGNN \cite{fu2020magnn}, MECCH \cite{fu2024mecch} and MF2Vec \cite{kim2024mf2vec} encode instances, but the result is written to the target node only. Our model keeps the original graph and writes the instance back to all of its nodes.

\paragraph{Sheaf neural networks.} Cellular sheaves attach a vector space and learned restriction maps to a graph, so that neighbours are compared after a change of coordinates \cite{hansen2020sheaf,bodnar2022neural,barbero2022sheaf}; extensions include cooperative and directional variants \cite{ribeiro2026cooperative,fiorini2025sheaves}. HetSheaf \cite{braithwaite2024heterogeneous} adapts the idea to heterogeneous graphs through edges only. Sheaf hypergraph networks \cite{duta2023sheafhyper} learn a map per node and hyperedge freely. Cell-complex and topological networks \cite{bodnar2021cwn,hajij2022tdl,hajij2025ctnn} pass messages over higher-dimensional cells, but without a sheaf whose face maps are fixed by the graph's own edge maps.

\section{Preliminaries}
\label{sec:background}

\paragraph{Heterogeneous graphs and meta-paths.} A heterogeneous graph is $G=(V,E,\phi,\psi)$ with node types $\phi:V\to\mathcal{T}$ and relation types $\psi:E\to\mathcal{R}$, where $|\mathcal{T}|>1$ or $|\mathcal{R}|>1$. We keep one edge per fact, not one per direction. A meta-path schema of length $\ell$ is a typed relation sequence $\pi: t_0\xrightarrow{r_1}t_1\cdots\xrightarrow{r_\ell}t_\ell$ \cite{sun2011pathsim}; a meta-path instance is a sequence $f=(v_0,e_1,v_1,\ldots,e_\ell,v_\ell)$ with $\phi(v_i)=t_i$ and $\psi(e_i)=r_i$. The schema is a type-level pattern, the instance is one occurrence of it. The nodes $v_1,\dots,v_{\ell-1}$ are intermediate; in Figure~\ref{fig:example}, PAP is a schema and P1--A7--P9 an instance with A7 intermediate.

\paragraph{Heterophily.} Labels are comparable only between target-type nodes, so heterophily is measured on the graphs that length-two schemas from the target type to itself induce. We report the $\mathcal{H}^2$ index of H2GB \cite{lin2025heterophily}, the mean over these graphs of $1-h_{\mathrm{adj}}$ with the class-imbalance-corrected adjusted homophily $h_{\mathrm{adj}}$ \cite{platonov2023characterizing}, and the MLH and MDE measures of Hetero$^2$Net \cite{li2025hetero2net}, each as its source defines it.

\paragraph{Sheaves on a graph and its incidence poset.} A \emph{cell} is the generic word for a node (0-cell), an edge (1-cell) or a face (2-cell). A cellular sheaf assigns every cell a vector space, its \emph{stalk} (we also call the vector it holds a stalk, and speak of node, edge and face stalks), and every incident pair $p\inc q$ a linear \emph{restriction map} $F_{p\inc q}$ from the stalk of $p$ to that of $q$, such that maps compose along chains: $F_{e\inc f}F_{v\inc e}=F_{v\inc f}$. On a graph, only the pairs $v\inc e$ (a node and an edge it lies on) exist, and the sheaf Laplacian measures $\sum_e\sum_{v\in e}\norm{F_{v\inc e}x_v-s_e}^2$ after eliminating the edge stalks $s_e$; with identity maps it is the graph Laplacian. We use the \emph{incidence poset} that adds faces: $v\inc e$ if $v$ lies on $e$, $e\inc f$ if $e$ is an edge of $f$, and $v\inc f$ if $v$ is a node of $f$.

%======================================================================
\section{Methodology}
\label{sec:method}

\subsection{Problem}
Given a heterogeneous graph with type-dependent node features $x^{\mathrm{raw}}_v\in\R^{d_{\phi(v)}}$ (types without features receive a learned vector per node) and no list of meta-path schemas, we learn a predictor $p_\theta(o;G)$ for an object $o$: a target node (classification) or a typed node pair (link prediction). The parameters minimise the average of a loss $\mathcal{L}$, cross-entropy or binary cross-entropy, over labelled training objects.

\subsection{Overview}
\begin{figure*}[t]
  \centering
  \includegraphics[width=\textwidth]{figure/h2snn_overview_v3b2.pdf}
  \caption{Overview of \method{}: (1) encode nodes, (2) learn restriction maps on edges, (3) sample faces with \wisp{}, (4) derive face maps from edge maps and weight each face, (5) combine node, edge and face stalks in one energy, (6) propagate by minimising the energy and predict.}
  \label{fig:pipeline}
\end{figure*}
\method{} has six steps (Figure~\ref{fig:pipeline}). Node features are encoded into a shared space. Each edge receives learned restriction maps. \wisp{} samples meta-path instances as faces; their maps follow from the edge maps, and each receives a weight. Node, edge and face stalks then interact through one energy, and the node representations are its minimiser. A prediction head reads each node's own encoding together with its minimiser.

\subsection{Encoding and Restriction Maps}
Nodes of different types must be compared by the same maps, so each node is first encoded into a $d$-dimensional vector, $h_v=\operatorname{GELU}(\operatorname{LN}(W_{\phi(v)}x^{\mathrm{raw}}_v+b_{\phi(v)})+q_{\phi(v)})$ with type-specific $W,b$ and a type embedding $q$. Two nodes with similar features can play different roles, so $c_v$ adds the mean of $h$ over the neighbours of $v$ under every relation and direction.

An edge $e=(u\xrightarrow{r}v)$ carries one learned restriction map per endpoint, which lets neighbours with different labels be consistent in different coordinates. Each map is orthogonal and depends on both endpoints, the relation and the role of the endpoint:
\begin{equation}
\begin{aligned}
  F_{v\inc e}&=R_{r,\rho}\,H_2(u)H_1(u),\\
  u&=\operatorname{MLP}([h_v,c_v,h_{v'},c_{v'},\zeta_r,\rho]),
\end{aligned}
  \label{eq:maps}
\end{equation}
where $v'$ is the other endpoint, $\zeta_r$ a relation embedding, $\rho$ the role, $R_{r,\rho}$ a learned rotation, and $H_1,H_2$ Householder reflections. Orthogonality keeps all compositions of maps orthogonal; the pair conditioning lets the identity of an intermediate node enter the geometry, which relation-only maps cannot do. Each edge also has a positive weight $a_e=c_e^2$.

\subsection{Faces: Meta-Path Instances with Forced Maps}
\label{sec:wisp}
\wisp{} supplies the faces without a meta-path list. From every target node it draws $r$ random walks on the original graph; a walk draws its length uniformly from $\{2,\ldots,\ell_{\max}\}$, chooses uniformly among relation--direction pairs that lead to an unvisited node, then uniformly among those neighbours, and is discarded at a dead end. A walk and its reversal give one instance. Each walk $\omega$ has a probability $p(\omega)$, and each instance receives the sampling weight $m_f=\sum_{\omega\to f}\frac1r\,p(\omega)/q_{\mathrm{mix}}(\omega)$ with $q_{\mathrm{mix}}$ the distribution the walks were drawn from, so that the weighted sum of instance terms is an unbiased estimate of the one defined by $p$ (Appendix~\ref{app:proofs}). Training draws new walks every epoch; evaluation uses one fixed draw, shared by all compared models.

A face is one instance $f=(v_0,e_1,\ldots,v_\ell)$. It is \emph{not} closed: there is no edge between $v_0$ and $v_\ell$. A face has a stalk $z_f$, and its restriction maps are not new parameters. We fix $F_{v_0\inc f}=I$ and compose the edge maps along the path,
\begin{equation}
\begin{aligned}
  F_{e_i\inc f}&=F_{v_{i-1}\inc f}\,F_{v_{i-1}\inc e_i}^{\transpose},\\
  F_{v_i\inc f}&=F_{e_i\inc f}\,F_{v_i\inc e_i}.
\end{aligned}
\label{eq:forced}
\end{equation}
Reading the first line, the map into the face from edge $e_i$ is the one that sends the stalk of $e_i$ back to the frame of $v_{i-1}$ and then forward to the face; the second line carries the frame on to $v_i$. By construction $F_{e_i\inc f}F_{v_{i-1}\inc e_i}=F_{v_{i-1}\inc f}$ and $F_{e_i\inc f}F_{v_i\inc e_i}=F_{v_i\inc f}$, so all three kinds of maps compose along the incidence poset.

\begin{proposition}[Forced maps]\label{prop:forced}
With Equation~\eqref{eq:forced}, the edge, face and node maps form one sheaf on the incidence poset, and the face terms of the energy below vanish at every global section of the graph's own sheaf. Freely learned face maps do not have this property and in general shrink the set of zero-energy configurations.
\end{proposition}
Hence adding faces never rejects a configuration that the edges already accept; it only adds constraints where the edges are silent.

\subsection{One Energy}
\label{sec:energy}
The model keeps three kinds of stalks: $x_v$ for every node, $s_e$ for every edge and $z_f$ for every face, collected in $Y=(x,s,z)$. The energy asks every comparable pair of cells to agree after their restriction maps:
\begin{equation}
\begin{aligned}
  E(Y)=&\tfrac12\sum_{e}a_e\!\!\sum_{v\in e}\!\norm{F_{v\inc e}x_v-s_e}^2\\
  +&\tfrac12\sum_{f}\beta_f\Bigl[\sum_{e\in f}\norm{F_{e\inc f}s_e-z_f}^2+\!\sum_{v\in f}\norm{F_{v\inc f}x_v-z_f}^2\Bigr].
\end{aligned}
\label{eq:energy}
\end{equation}
The first sum compares each node with its edges; the second compares each face with its edges and with every one of its nodes, intermediate ones included. The weight $\beta_f=w_fm_f$ combines the sampling weight with a learned score $w_f=\exp g(\xi_f)$, where $g$ reads the encodings $h$ of all nodes of the instance and its relation embeddings, so two instances of one schema can differ in influence. For the instance P1--A7--P9, the energy has four node--edge pairs, two edge--face pairs and three node--face pairs; the last include A7, so its stalk is pulled toward the common value $z_f$ together with P1 and P9, and passes the result to P4 through its own edges.

\begin{proposition}[Reach]\label{prop:reach}
Let $L=\nabla^2E$. For every node $v$ of face $f$, the $(x_v,z_f)$ block of $L$ is $-\beta_fF_{v\inc f}^{\transpose}$, which is nonzero. In an operator that couples a face only to its edges, this block is zero, and a node's stalk can reach the face only through edges.
\end{proposition}

\subsection{Propagation as a Minimiser and Prediction}
\label{sec:prop}
Let $D$ be the block diagonal of $L$ and $\hat L=D^{-1/2}LD^{-1/2}$. Node, edge and face stalks are placed on the same scale, so that influence does not grow with degree. Only node stalks have an input, namely the encoding $h$; edge and face stalks are free. The propagated stalks solve
\begin{equation}
  Y^\star=\arg\min_Y\ \tfrac12\norm{x-h}^2+\tfrac{\tau}{2}\,Y^{\transpose}\hat LY,\qquad \tau=e^{\theta},
  \label{eq:solve}
\end{equation}
that is, each node stays close to its encoding while the energy is small. $\tau$ is learned, with a larger learning rate than other parameters. The prediction is $\hat y_v=\operatorname{MLP}([x^\star_v\concat h_v])$, so a node's own encoding is not mixed with disagreeing neighbours before the nonlinear head \cite{zhu2020beyond}; link prediction scores a pair of such vectors with a relation-specific bilinear decoder.

\begin{proposition}[Unique solution]\label{prop:unique}
For every parameter value, Equation~\eqref{eq:solve} has a unique minimiser, $x^\star=(I+\tau\hat S)^{-1}h$ where $\hat S$ is the Schur complement of $\hat L$ onto the node block, and $\operatorname{cond}(I+\tau\hat S)\le1+\tau$.
\end{proposition}
We solve the system by preconditioned conjugate gradients and differentiate through the solution implicitly, so no unrolled layers are stored. Two facts keep the claim modest. The node-level map $h\mapsto x^\star$ is a resolvent, the same form as personalised PageRank propagation \cite{gasteiger2019appnp}; what is new is the operator inside it. And since edges and faces have no input features, their stalks carry no information beyond what the nodes supply; they change the geometry in which nodes interact.

\paragraph{Why this addresses heterophily.} Heterophily means that adjacent nodes should \emph{not} agree in the raw space. Edge maps let neighbours agree in another coordinate system, and a global section can place different classes in different directions of the stalk. Faces add a comparison that no edge expresses: P1 and P9 are compared \emph{through} A7, with A7's own maps, and the comparison is also available to A7 and to every node reached from A7. Existence of such class-separating sections under heterophily holds for the family of maps used here; it does not imply that training finds them.

\subsection{Special Cases and Cost}
\label{sec:cases}
The model contains the following configurations, which we compare on identical sampled instances. \emph{No faces}: $\beta_f=0$, an edge-sheaf model. \emph{Endpoint edge}: each face is replaced by one edge between $v_0$ and $v_\ell$ with the composed map, which is what remains after eliminating intermediate nodes. \emph{Face as feature}: the face stalk is computed from the instance and added to its target node only, with no coupling back; with relation-only maps it equals the instance encoder of MAGNN, as we checked numerically. \emph{No node--face} and \emph{No edge--face}: the corresponding terms of Equation~\eqref{eq:energy} are dropped. \emph{Free face maps}: face maps are learned. The weight ladder restricts $g$ to a constant (\emph{Uniform}), to the relation sequence (\emph{Schema}) or to the endpoints and relation sequence (\emph{Endpoints}); the full model additionally sees the intermediate nodes.

Edges that lie on at most one face are eliminated exactly, so the cost depends on the edges shared by several faces. On ACM, DBLP and IMDB with $d=32$, the solver converged in 9, 11 and 10 iterations at tolerance $10^{-4}$, and peak memory was measured on one RTX 4060 (8 GB) card; wall-clock times against baselines are in Table~\ref{tab:cost}. The cost does not depend on the number of schemas.

%======================================================================
\section{Experiments}
\label{sec:experiments}

\subsection{Experiment Settings}
We evaluate on node classification (main task) and link prediction; graph classification checks that the model covers the task and is reported in Appendix~\ref{app:data}.

\paragraph{Datasets.} For node classification we use four HGB benchmarks \cite{lv2021hgb}, ACM, DBLP, IMDB and Freebase (24\%/6\%/70\% splits), and four H2GB benchmarks \cite{lin2025heterophily}: ogbn-mag (time split 85/9/6), mag-year (random 50/25/25), IEEE-CIS-G (time split 80/10/10) and H-Pokec (random 50/25/25). They cover a wide range of heterophily (Table~\ref{tab:hetero}). For link prediction we use eight tasks: ACM paper--author, DBLP author--paper, IMDB movie--actor, Freebase book--book, PubMed disease--disease, LastFM user--artist, ogbn-mag paper--paper and H-Pokec user--user. Held-out facts are 10\% of the target relation (81/9/10 split by fact; Freebase by unordered pair) and are removed in both directions from the message graph and from every sampled instance; features computed from the target relation or describing the other endpoint are removed, and a leakage audit is run.

\paragraph{Baselines.} We compare with homogeneous homophilic (GCN \cite{kipf2017gcn}, GAT \cite{velickovic2018gat}, ChebNetII \cite{he2022chebnetii}), homogeneous heterophilic (NSD \cite{bodnar2022neural}, CSNN \cite{ribeiro2026cooperative}), heterogeneous homophilic (R-GCN \cite{schlichtkrull2018rgcn}, HAN \cite{wang2019han}, MAGNN \cite{fu2020magnn}, SeHGNN \cite{yang2023sehgnn}, HiNormer \cite{mao2023hinormer}, PSHGCN \cite{he2024spectral}) and heterogeneous heterophilic models (H2G-former \cite{lin2025heterophily}, Hetero$^2$Net \cite{li2025hetero2net}, H2SGNN \cite{lu2026addressing}, HetSheaf \cite{braithwaite2024heterogeneous}, LatGRL \cite{shen2025heterophily}). Following HGB, the groups are named Hom./hom., Hom./het., Het./hom. and Het./het. by whether the graph is modeled as homogeneous or heterogeneous and the method designed for homophilic or heterophilic settings. Link prediction uses a subset covering all groups.

\paragraph{Protocol.} Node classification is scored by Macro-F1 and Micro-F1, with Macro-F1 averaged over the classes present in each split's true labels; IMDB is multi-label (binary cross-entropy per genre, threshold at logit zero, 4,573 labelled movies). Model selection uses validation Macro-F1. For link prediction we follow the HGB protocol: each test positive is paired with one negative whose tail is a two-hop neighbour of the head within the target type; validation uses random tail negatives; we report ROC-AUC and MRR and select on validation ROC-AUC. All methods, baselines included (HetSheaf through an adapter), are rerun under one protocol with their original feature preprocessing, plain Bayesian search with 100 trials on seed 0 and then seeds 0--4. For \method{} the search covers $d$, learning rate, weight decay, and encoder and classifier dropout. We report mean $\pm$ standard deviation over five seeds and test each comparison with Welch's $t$-test, reporting raw and adjusted $p$-values with the adjustment chosen after the results. \wisp{} settings are fixed before any run and never tuned: $\ell_{\max}=4$, with $r$ calibrated so that the number of node--face incidences matches that of the literature meta-paths on the same dataset, and a uniform proposal.

\begin{table}[t]
  \centering
  \caption{Heterophily of the node-classification datasets, computed with the reference implementations; higher $\mathcal{H}^2$ and MLH mean more heterophilic. Reproduction target: Hetero$^2$Net's published MLH and MDE.}
  \label{tab:hetero}
  \footnotesize
  \setlength{\tabcolsep}{2.4pt}
  \begin{tabular}{lcccccccc}
    \toprule
    & ACM & DBLP & IMDB & Free. & mag & m-yr & IEEE & Pok. \\
    \midrule
    $\mathcal{H}^2$ & \res & \res & \res & \res & \res & \res & \res & \res \\
    MLH & \res & \res & \res & \res & \res & \res & \res & \res \\
    MDE & \res & \res & \res & \res & \res & \res & \res & \res \\
    \bottomrule
  \end{tabular}
\end{table}

\subsection{Results}
\label{sec:results}
Table~\ref{tab:node-cls} reports node classification and Table~\ref{tab:link-pred} link prediction. \res{}

\begin{table*}[t]
  \centering
  \caption{Node classification, Macro-F1, mean $\pm$ std over five seeds (Micro-F1 in Appendix~\ref{app:data}). Datasets are ordered by $\mathcal{H}^2$; best values are bold.}
  \label{tab:node-cls}
  \footnotesize
  \setlength{\tabcolsep}{4pt}
  \begin{tabular}{llcccccccc}
    \toprule
    \textbf{Group} & \textbf{Model} & ACM & DBLP & IMDB & Freebase & ogbn-mag & mag-year & IEEE-CIS-G & H-Pokec \\
    \midrule
    \multirow{3}{*}{Hom./hom.} & GCN & \res & \res & \res & \res & \res & \res & \res & \res \\
    & GAT & \res & \res & \res & \res & \res & \res & \res & \res \\
    & ChebNetII & \res & \res & \res & \res & \res & \res & \res & \res \\
    \midrule
    \multirow{2}{*}{Hom./het.} & NSD & \res & \res & \res & \res & \res & \res & \res & \res \\
    & CSNN & \res & \res & \res & \res & \res & \res & \res & \res \\
    \midrule
    \multirow{6}{*}{Het./hom.} & R-GCN & \res & \res & \res & \res & \res & \res & \res & \res \\
    & HAN & \res & \res & \res & \res & \res & \res & \res & \res \\
    & MAGNN & \res & \res & \res & \res & \res & \res & \res & \res \\
    & SeHGNN & \res & \res & \res & \res & \res & \res & \res & \res \\
    & HiNormer & \res & \res & \res & \res & \res & \res & \res & \res \\
    & PSHGCN & \res & \res & \res & \res & \res & \res & \res & \res \\
    \midrule
    \multirow{5}{*}{Het./het.} & H2G-former & \res & \res & \res & \res & \res & \res & \res & \res \\
    & Hetero$^2$Net & \res & \res & \res & \res & \res & \res & \res & \res \\
    & H2SGNN & \res & \res & \res & \res & \res & \res & \res & \res \\
    & HetSheaf & \res & \res & \res & \res & \res & \res & \res & \res \\
    & LatGRL & \res & \res & \res & \res & \res & \res & \res & \res \\
    \midrule
    Ours & \method{} & \res & \res & \res & \res & \res & \res & \res & \res \\
    \bottomrule
  \end{tabular}
\end{table*}

\begin{table*}[t]
  \centering
  \caption{Link prediction, ROC-AUC (MRR in Appendix~\ref{app:data}), mean over five seeds.}
  \label{tab:link-pred}
  \footnotesize
  \setlength{\tabcolsep}{4pt}
  \begin{tabular}{lcccccccc}
    \toprule
    \textbf{Model} & ACM & DBLP & IMDB & Freebase & PubMed & LastFM & ogbn-mag & H-Pokec \\
    \midrule
    GCN & \res & \res & \res & \res & \res & \res & \res & \res \\
    NSD & \res & \res & \res & \res & \res & \res & \res & \res \\
    R-GCN & \res & \res & \res & \res & \res & \res & \res & \res \\
    MAGNN & \res & \res & \res & \res & \res & \res & \res & \res \\
    HetSheaf & \res & \res & \res & \res & \res & \res & \res & \res \\
    H2G-former & \res & \res & \res & \res & \res & \res & \res & \res \\
    \midrule
    \method{} & \res & \res & \res & \res & \res & \res & \res & \res \\
    \bottomrule
  \end{tabular}
\end{table*}

%======================================================================
\section{Discussion}
\label{sec:discussion}

\subsection{What the Faces Contribute}
To separate the contribution of faces from that of weighting and capacity, the configurations of Section~\ref{sec:cases} are compared on identical sampled instances, the same evaluation draw and parameter counts within 5\% (Table~\ref{tab:ablation}). \emph{Endpoint edge} tests whether the intermediate nodes matter, \emph{Face as feature} whether coupling back to all nodes matters beyond a MAGNN-type encoding, and \emph{Free face maps} whether forcing the maps matters. The criterion is fixed before running: if the full model is not significantly better than the \emph{Endpoint} weight variant on at least two of ACM, DBLP and IMDB, we do not claim that intermediate nodes contribute through the weights. \res{}

\begin{table}[t]
  \centering
  \caption{Controlled variants on identical sampled instances (Macro-F1, mean over five seeds).}
  \label{tab:ablation}
  \footnotesize
  \setlength{\tabcolsep}{3pt}
  \begin{tabular}{lcccc}
    \toprule
    \textbf{Variant} & ACM & DBLP & IMDB & Freebase \\
    \midrule
    \method{} (full) & \res & \res & \res & \res \\
    No faces & \res & \res & \res & \res \\
    Endpoint edge & \res & \res & \res & \res \\
    Face as feature & \res & \res & \res & \res \\
    No node--face term & \res & \res & \res & \res \\
    No edge--face term & \res & \res & \res & \res \\
    Free face maps & \res & \res & \res & \res \\
    Uniform / Schema / Endpoints weights & \res & \res & \res & \res \\
    \bottomrule
  \end{tabular}
\end{table}

\subsection{Heterophily and Cost}
Reviewers of an earlier version asked for heterophily to be established and analysed, and for cost to be reported against graph size. We therefore report accuracy of the full model against the baselines by $\mathcal{H}^2$ and by node degree, and time and memory per epoch with graph size (Table~\ref{tab:cost}). \res{}

\begin{table}[t]
  \centering
  \caption{Cost per epoch with graph size.}
  \label{tab:cost}
  \footnotesize
  \setlength{\tabcolsep}{4pt}
  \begin{tabular}{lccc}
    \toprule
    & ACM & DBLP & IMDB \\
    \midrule
    Time per epoch (s) & \res & \res & \res \\
    Peak memory (MB) & \res & \res & \res \\
    \bottomrule
  \end{tabular}
\end{table}

\subsection{Limitations}
Edges, faces and the energy are symmetric, so the strength of influence between two nodes is the same in both directions; the energy is quadratic and the propagation linear, so no band-pass response is built in; faces are simple walks, so branching meta-structures are not represented; and the weights scale a face but cannot change its maps, which the edges fix. Existence results do not imply that training finds a class-separating section.

%======================================================================
\section{Conclusion}
We asked whether heterophily can be handled on the original heterogeneous graph by making meta-path instances part of its structure. \method{} does so by treating each sampled instance as a face of the graph's own sheaf, with maps forced by its edges, coupled to every node of the instance and solved with the nodes and edges in one energy with a unique minimiser. \res{}

\bibliographystyle{named}
\bibliography{references}

%======================================================================
\appendix
\section{Proofs}
\label{app:proofs}
Every edge map is orthogonal, so the maps of Equation~\eqref{eq:forced} are orthogonal and $F^{\transpose}F=I$.

\emph{Unbiased sampling.} Walks are drawn independently, $r$ per anchor, from $q_{\mathrm{mix}}$. Then $\mathbb{E}[\sum_fm_fL_f]=\sum_a r\sum_\omega q_{\mathrm{mix}}(\omega)\frac1r\frac{p(\omega)}{q_{\mathrm{mix}}(\omega)}L_{f(\omega)}=\sum_a\sum_\omega p(\omega)L_{f(\omega)}$.

\emph{Proposition~\ref{prop:forced}.} Composition is verified in Equation~\eqref{eq:forced}. If $F_{v\inc e}x_v=s_e$ for all incident pairs, set $z_f=x_{v_0}$ and observe $F_{e_i\inc f}s_{e_i}=F_{v_{i-1}\inc f}x_{v_{i-1}}=z_f$ by induction and $F_{v_i\inc f}x_{v_i}=F_{e_i\inc f}s_{e_i}=z_f$, so every face term vanishes.

\emph{Proposition~\ref{prop:reach}.} Differentiate Equation~\eqref{eq:energy} with respect to $x_v$ and $z_f$: the only term containing both is $\frac12\beta_f\norm{F_{v\inc f}x_v-z_f}^2$, giving $-\beta_fF_{v\inc f}^{\transpose}$, nonzero because $F_{v\inc f}$ is orthogonal and $\beta_f>0$.

\emph{Proposition~\ref{prop:unique}.} $L_{xx}$ is block diagonal with $D_vI$ because there are no node--node terms, so $\hat L_{xx}=I$. The block of edge and face stalks, $\hat L_{ww}$, is positive definite: its quadratic form is zero only if every $s_e$ and then every $z_f$ is zero, since $a_e,\beta_f>0$. Hence the Schur complement $\hat S=I-\hat L_{xw}\hat L_{ww}^{-1}\hat L_{wx}$ satisfies $0\preceq\hat S\preceq I$, the objective is strictly convex, $x^\star=(I+\tau\hat S)^{-1}h$, and $\operatorname{cond}(I+\tau\hat S)\le1+\tau$.

\section{Datasets, Settings and Full Results}
\label{app:data}
This appendix reports dataset statistics, the Micro-F1 and MRR tables, search spaces, graph classification and further results. Graph classification follows HetSheaf's protocol (TU MUTAG, PROTEINS, ENZYMES, NCI1 with node types taken from the label one-hot, constant features, 80/10/10 split, seeds 0--4); because node types are derived from labels, it checks generalisability only. \res{}

\end{document}
